\section{Future Work}

A next step is to introduce finer-grained intermediate architectures between our simple baselines and the full gateway setup. 
Our current comparison involves a relatively large architectural step, which makes it harder to attribute the observed gains to any single design choice. 
Future work should therefore minimize framework-specific differences between the two architecture families and introduce incremental variants that isolate individual success drivers more precisely.

At the same time, the gateway requires deeper investigation of its security properties before it can be considered for industry deployment. 
As protocol-based agent systems grow in scope, authentication, attack resistance, and trust boundaries become critical concerns that should be addressed as part of future gateway research.

Our evaluation also revealed the lack of a standardized benchmark for multi-agent systems that operate across multiple protocols. 
Existing benchmarks focus primarily on MCP tool retrieval~\cite{mo2026livemcpbench, luo2025mcpuniverse}, which does not capture the coordination and protocol-switching demands of gateway-based setups. 
We therefore had to design our own task environment to evaluate the different pipeline architectures. 
A dedicated multi-protocol benchmark would benefit future research in this area considerably.

On the system side, the gateway could be extended with more specialized retrieval methods.
Prior work already explores promising tool-retrieval strategies~\cite{gaurav2025dynamicreact, fei2025mcpzero, gan2025ragmcp, lumer2025tooltoagent, lumer2025scalemcp}.
Recent work also shows that specialized retrieval models can improve tool retrieval further~\cite{lu2025toolsunderdocumented}.

Finally, our token consumption metric currently suffers from excessive upstream token usage caused by results being passed back through the whole agent chain without contributing additional value. 
Future work should investigate strategies to reduce this overhead, as it inflates measured costs in multi-agent pipelines like ours and limits the scalability of deeper delegation chains.
